\section{Turn the message into complete blocks}
A compression call consumes exactly 64 bytes. A message whose length is not a multiple
of 64 leaves an incomplete block. Even a message that fills its last block exactly
needs a final block that records its end. SHA-256 handles both cases by adding padding
once, after the last message byte.

For a byte-oriented message, append the byte \texttt{80}. This contributes a one bit
followed by seven zero bits. Append zero bytes until eight bytes remain at the end of
a 64-byte block. Those last eight bytes hold the original message length in bits,
most significant byte first. The padding bytes themselves do not increase the recorded
message length.

For \texttt{abc}, the original length is three bytes, or 24 bits. The first four
bytes of the padded block are \texttt{61 62 63 80}. Bytes 4 through 55 are zero.
Bytes 56 through 63 are \texttt{00 00 00 00 00 00 00 18}. In this paragraph the
byte indices are decimal and the byte values are hexadecimal. The \texttt{18}
therefore represents decimal 24.

\begin{figure}[htbp]
  \centering
  \begin{tikzpicture}[alt={The padded abc block contains message bytes 61, 62 and 63 at positions zero through two, the byte 80 at position three, zero bytes through position 55, and an eight-byte big-endian length equal to hexadecimal 18 at positions 56 through 63.},x=1cm,y=0.8cm,font=\small]
    \draw (0,0) rectangle (2,-1);\draw (2,0) rectangle (3.4,-1);
    \draw (3.4,0) rectangle (7,-1);\draw (7,0) rectangle (10,-1);
    \node at (1,-0.5) {\texttt{61 62 63}};\node at (2.7,-0.5) {\texttt{80}};
    \node at (5.2,-0.5) {52 zero bytes};\node at (8.5,-0.5) {\texttt{00 \ldots\ 00 18}};
    \node at (1,0.5) {Message};\node at (2.7,0.5) {Marker};
    \node at (5.2,0.5) {Fill};\node at (8.5,0.5) {64-bit bit count};
    \node at (1,-1.5) {0--2};\node at (2.7,-1.5) {3};
    \node at (5.2,-1.5) {4--55};\node at (8.5,-1.5) {56--63};
  \end{tikzpicture}
  \caption{One padded block for the three-byte message \texttt{abc}. Positions are decimal.}
\end{figure}

The boundary is easiest to understand by counting free bytes. An incomplete block with
55 message bytes has nine free bytes: one for \texttt{80} and eight for the length.
With 56 message bytes it has only eight free bytes. Write the marker, fill the rest
of that block with zeros, compress it, and put the length at the end of another
zero-filled block. A 64-byte message has already been compressed by update; its
empty remaining buffer receives a complete padding block.

\begin{table}[htbp]
  \centering\caption{Padding around a block boundary. All counts are decimal.}
  \tagpdfsetup{table/header-rows={1}}
  \begin{tabular}{rrrr}
    \toprule
    Message bytes & Full blocks before final & Buffered bytes & Final blocks \\
    \midrule
    0             & 0                        & 0              & 1            \\
    3             & 0                        & 3              & 1            \\
    55            & 0                        & 55             & 1            \\
    56            & 0                        & 56             & 2            \\
    63            & 0                        & 63             & 2            \\
    64            & 1                        & 0              & 1            \\
    65            & 1                        & 1              & 1            \\
    120           & 1                        & 56             & 2            \\
    \bottomrule
  \end{tabular}
\end{table}
